%%*************************************************************************
%%
%% Křívky učení v raném stadia učení, anebo setkání člověka s cizím nastrojem.
%% V0.0
%% 2012/11/28
%% Timothy Hobbs
%% timothyhobbs@seznam.cz
%%
%%*************************************************************************
%%
%% Legal Notice:
%%
%% This code is offered as-is without any warranty either expressed or
%% implied; without even the implied warranty of MERCHANTABILITY or
%% FITNESS FOR A PARTICULAR PURPOSE! 
%% User assumes all risk.
%%
%% This work by Timothy Hobbs is licensed under a
%% Creative Commons Attribution-NonCommercial-ShareAlike 3.0 Unported License.
%% http://creativecommons.org/licenses/by-nc-sa/3.0/
%%
%% Loriem ipsum was from
%% https://github.com/pborky/documents/tree/master/DIP
%% by Peter Boraros
%%*************************************************************************
\documentclass[a4paper,12pt,twoside,czech,cleardoublepage=empty]{scrreprt}
\usepackage[T1]{fontenc}
\usepackage[utf8]{inputenc}
\usepackage{textcomp}
\usepackage[puttinydots,useemptybox,8dots]{braille}
\usepackage[czech]{babel}%počeštění názvů (Obsah, Kapitola, Literatura atp.)
%<------------------------------ headers ------------------------------------>
\newcommand{\SubjectCs}{Bakalařská Prace}
\newcommand{\TitleCs}{Křívky učení v raném stadia učení, anebo setkání člověka s cizím nastrojem.}
\newcommand{\TitleEn}{Learning curves in early stages of skill aquisition.}
\newcommand{\KeywordsEn}{Learning Curves, Braille, Reading Speed}
\newcommand{\KeywordsCs}{Křívky učení, Braillové pismo, rychlost čtení}
\newcommand{\AuthorName}{Timothy Vladimir Hobbs}
\newcommand{\SupervisorName}{Ing. Karolina Duschinská, Ph.D.}
\newcommand{\SupervisorEmail}{karolina.duschinska@pedf.cuni.cz}
\newcommand{\IssuedIn}{Praha}
\newcommand{\DepartmentEn}{Department of Pedagogy}
\newcommand{\DepartmentCs}{Katedra pedagogiky}
\newcommand{\Thanks}{Poděkování,

Děkují Samuel Thibault, který podílel na navrhování nízko úrovňový protokol. Dale, David Mielke, který vyvijel ovladač těchto zařízení. Členové brmlab.cz, které pomohli na technické straně navrhování. Christopher Olah, který vyvijel knihovnu ImplicitCad, s kterou jsem navrhl plastové díly pro toho zařízení.  Členové freenode síť IRC, zejmená členu kanalů \#\#c, \#arduino, a \#haskell pro jejích podporu na softwarové straně tehle studie.  Knihovna a Tiskárna Pro Nevidomé K.E. Macana, která poskytla pro mně bezplatnou tiskářskou službu.  Gimnázium pro zrakově postižené a Střední odborná škola pro zrakově postižené v Praze v Praze kde tento výzkum proběhel. Vedouci mého bakalářské prace, Ing. Karolina Duschinská, Ph.D..}

%<------------------------------         ------------------------------------>

%\usepackage{array}
%\usepackage{longtable}
%\usepackage{varioref}
%\usepackage{wrapfig}
%\usepackage{fancybox}
%\usepackage{framed}
\usepackage{url}
%\usepackage{amssymb,amsmath}

%\usepackage{graphicx}

\usepackage{ifpdf}
\ifpdf
\usepackage[pdftex]{graphicx}
\graphicspath{{./img/}}
\DeclareGraphicsExtensions{.pdf}
\else
\usepackage[dvips]{graphicx}
\graphicspath{{./img/}}
\DeclareGraphicsExtensions{.eps}
\fi

\usepackage{caption}
\usepackage{subcaption}


\makeatletter

%%%%%%%%%%%%%%%%%%%%%%%%%%%%%% LyX specific LaTeX commands.
%\providecommand{\LyX}{L\kern-.1667em\lower.25em\hbox{Y}\kern-.125emX\@}
\newcommand{\lyxline}[1][1pt]{%
  \par\noindent%
  \rule[.5ex]{\linewidth}{#1}\par}
\newcommand{\noun}[1]{\textsc{#1}}
% Special footnote code from the package 'stblftnt.sty'
% Author: Robin Fairbairns -- Last revised Dec 13 1996
%\let\SF@@footnote\footnote
%\def\footnote{\ifx\protect\@typeset@protect
%    \expandafter\SF@@footnote
%  \else
%    \expandafter\SF@gobble@opt
%  \fi
%}
%\expandafter\def\csname SF@gobble@opt \endcsname{\@ifnextchar[%]
%  \SF@gobble@twobracket
%  \@gobble
%}
%\edef\SF@gobble@opt{\noexpand\protect
%  \expandafter\noexpand\csname SF@gobble@opt \endcsname}
%\def\SF@gobble@twobracket[#1]#2{}
%% Because html converters don't know tabularnewline
%\providecommand{\tabularnewline}{\\}

%%%%%%%%%%%%%%%%%%%%%%%%%%%%%% Textclass specific LaTeX commands.
%\newenvironment{lyxcode}
%{\begin{list}{}{
%\setlength{\rightmargin}{\leftmargin}
%\setlength{\listparindent}{0pt}% needed for AMS classes
%\raggedright
%\setlength{\itemsep}{0pt}
%\setlength{\parsep}{0pt}
%\normalfont\ttfamily}%
% \item[]}
%{\end{list}}

%%%%%%%%%%%%%%%%%%%%%%%%%%%%%% User specified LaTeX commands.
%<-------------------------------společná nastavení------------------------------>
%\usepackage[]{hyperref} %odkazy v  pdf jsou klikací s barevnými rámečky
%\usepackage[urlcolor=blue]{hyperref}
%\usepackage[colorlinks=true,urlcolor=blue]{hyperref}
\usepackage{natbib} %balíček pro citace literatury  
%\usepackage{hypernat}%interakce mezi hyperref a natbib
%\newcommand{\BibTeX}{{\sc Bib}\TeX}%BibTeX logo
%\hypersetup{   % Nastavení polí PDF dokumentu 
%pdftitle={\TitleCs},%
%pdfauthor={\AuthorName},%
%pdfsubject={\SubjectCs},%
%pdfkeywords={\KeywordsCs}%
%}
\usepackage{multicol}

%<-----------------------------volání stylů----------------------------------------->
% (znak % je označení komentáře: co je za ním, není aktivní)
%<------------------------------------písmo----------------------------------------->
%\usepackage{packages/bc-latinmodern}
\usepackage{packages/bc-times}
%\usepackage{packages/bc-palatino}
%\usepackage{packages/bc-iwona}
%\usepackage{packages/bc-helvetika}


%<------------------------------záhlaví stránek------------------------------------>
%\usepackage{packages/bc-headings}
%\usepackage{packages/bc-fancyhdr}

%<------------------------------hlavičky kapitol------------------------------------>
%\usepackage{packages/bc-neueskapitel}
%\usepackage{packages/bc-fancychap}

\makeatother

\begin{document}
~\thispagestyle{empty}\vfill{}
%Tato stránka je tzv. protititul a je graficky součástí titulní stránky.
%Nechte ji prázdnou, nebo na ni umístěte vhodnou fotografii či ilustraci.

%\csprimeson             % `` '' se budou sazet jako ceske
\clearpage{}~\thispagestyle{empty}\begin{center}\pagenumbering{roman}\vspace{10mm}


\textsf{\textsc{\noun{\LARGE Universita Karlova v Praze}}}\\
\vspace{0.5em}
\textsf{\textsc{\noun{\LARGE Pedagogická Fakulta}}}\\
\vspace*{1em}
\textsf{\textsc{\noun{\Large \DepartmentCs}}}\vspace{15mm}


%\includegraphics[width=0.3\textwidth]{uk}
\vspace{15mm}

\textsf{\textsc{\noun{\huge \SubjectCs}}}{\huge \par}

\vspace{15mm}


\textsf{\textsc{\noun{\LARGE \TitleCs}}}{\LARGE \par}


\vspace{10mm}


\end{center} 

\vspace*{\fill}


\vspace{10mm}


\begin{description}
\item [{{\large Autor:}}] \noindent \textsf{\large \AuthorName}{\large \par}
\item [{{\large Vedoucí}}] \noindent \textsf{\large \SupervisorName}%
{\large \hfill{}}\textsf{\large \IssuedIn, \the\year}{\large{}
}{\large \par}
\end{description}

\clearpage{}~\thispagestyle{empty}

{\small %\setcounter{page}{3} % nastavení číslování stránek
\ }{\small \par}

\noindent {\small \vfill{}
 % nastavuje dynamické umístění následujícího textu do spodní části stránky
~}{\small \par}

\noindent {\small Prohlašuji, že jsem svou diplomovou práci napsal(a)
samostatně a výhradně s použitím citovaných pramenů. Souhlasím se
zapůjčováním práce a jejím zveřejňováním. Reprodukce pod licence Creative Commons Zero je povolený. }{\small \par}

{\small \bigskip{}
}\noindent {\small{} V Praze dne \today\hspace{\fill} \AuthorName}\\
{\small{} % doplňte patřičné datum, jméno a příjmení
}{\small \par}

{\small %%%   Výtisk pak na tomto míste nezapomeňte PODEPSAT!
%%%                                         *********
}{\small \par}


\clearpage{}~\thispagestyle{empty}

~\thispagestyle{empty}{\small ~\vfill{}
}{\small \par}

\noindent {\small \Thanks \newpage{}}{\small \par}

\clearpage{}

{\small \thispagestyle{plain}}{\small \par}

\noindent {\small ~\vfill{}
}{\small \par}

\addcontentsline{toc}{chapter}{Abstract} 

\begin{description}
\item [{{\small Název~práce:}}] \noindent {\small \TitleCs}{\small \par}
\item [{{\small Autor:}}] \noindent {\small \AuthorName}{\small \par}
\item [{{\small Katedra~(ústav):}}] \noindent {\small \DepartmentCs}{\small \par}
\item [{{\small Vedoucí~diplomové~práce:}}] \noindent {\small \SupervisorName}
\item [{{\small e-mail~vedoucího:}}] \noindent {\small \SupervisorEmail}\\
{\small \par}
\item [{{\small Abstrakt}}] \noindent {\small Tato práce se zabývá prvními dojmy, adaptací a učením s podporou nového nástroje. V této studii jsem zkoumal, jak lidé se zrakovým postižením reagují na nový druh braillského řádku typu Fast Character Display (FCHAD) a jak se ho naučí používat. Pedagogické prostředí této studie je kontrolované a učení je založené na drilování. Zkoumání se zvláště zaměří na otázku adaptace během drilování. Studie se podobá psychologickým výzkumům křivek učení s tím rozdílem, že pedagogické působení je plnohodnotnou součástí studie.
 }{\small \par}
\item [{{\small Klíčová~slova:}}] \noindent {\small \KeywordsCs}\\
{\small \lyxline{\small}}{\small \par}
\item [{{\small Title:}}] \noindent {\small \TitleEn}{\small \par}
\item [{{\small Author:}}] \noindent {\small \AuthorName}{\small \par}
\item [{{\small Department:}}] \noindent {\small \DepartmentEn}{\small \par}
\item [{{\small Supervisor:}}] \noindent {\small \SupervisorName}{\small \par}
\item [{{\small Supervisor's~e-mail~address:}}] \noindent {\small \SupervisorEmail}\\
{\small \par}
\item [{{\small Abstract}}] \noindent {\small An examination of first impressions, adaptation and learning over a short period with novel material.
 In this study, I have examined how subjects react to, and learn to use, an FCHAD braille display.  The pedagogic atmosphere of this study is strict and drill based.  Particular focus is placed on how students adapt during the course of a drill.  The study was performed similarly to many learning curve studies, yet differs in that a teaching element was explicitly included rather than explicitly excluded.
 }{\small \par}
\item [{{\small Keywords:}}] \noindent {\small \KeywordsEn}{\small \par}
\end{description}
%\thispagestyle{empty}~{\small \addcontentsline{toc}{chapter}{Assignment of the Work} }{\small \par}


\thispagestyle{empty}{\small \tableofcontents{}% vkládá automaticky generovaný obsah dokumentu
\cleardoublepage{}}{\small \par}
\pagenumbering{arabic}%start arabic pagination from 1 
\renewcommand{\baselinestretch}{1.5}

\chapter{Úvod}

\section{Nevidomost ve světě a význam Braillova písma}

Nevidomost je fenomén rozvojových zemí a starších lidí. Na světě žije 39 milionů nevidomých. Devadesát procent z nich bydlí v rozvojových zemích a 65\% je starších padesáti let \citep{whodata}.

Gramotnost nevidomých je schopnost číst Braillovo písmo.  Nevidomí sami říkají, že čtou rukama.

Ve vyspělých zemích gramotnost souvisí se zaměstnaností. Podle statistiky zveřejněné ve Spojených státech je 90\% gramotných nevidomých zaměstnáno oproti pouhým 30 procentům negramotných nevidomých.

Pro Českou republiku ani rozvojový svět jsem podobná čísla nenašel.  Do jisté míry jsou statistiky zkreslené.  Ve Spojených státech gramotnost značně poklesá.  V roce 1960 se učilo číst rukama padesát procent nevidomých amerických žáků. Dnes je to jen dvanáct procent.  Selhání školského systému není jediným důvodem. Dříve mnoho lidí ztratilo zrak kvůli infekčním nemocem, což se dnes už téměř nestává. Dále už umíme léčit bílý zákal a ve vyspělých zemích neexistuje nevidomost jím způsobená.  Je tedy méně nevidomých, ale zato jsou více nemocní.  Vysoké procento nevidomých ve vyspělém světě trpí další poruchou, která je příčinou jejich negramotnosti.  Mnoho nevidomých se narodilo předčasně nebo mají perinatální poruchu  \citep{perkins,whozakal,whodata}.  Dva z mých sedmi nevidomých účastníků trpěli poruchou sluchu.

Dalším důvodem nízké gramotnosti nevidomých je omezená dostupnost materiálů v Braillově písmu.  V Knihovně a tiskárně pro nevidomé K.E. Macana je 5693 titulů tištěných v Braillově písmu \citep{biblio}. V Městské knihovně v Praze je naproti tomu 387625 knih pro vidomé  \citep{mlp}.

Typové složení nevidomosti v rozvojových zemích je výrazně jiné než ve vyspělém světě. Nevidomost ve světě je z osmdesáti procent léčitelná nebo se jí dokonce dalo vyhnout prevencí \citep{whodata}.  Celých 50\% nevidomosti je totiž způsobeno bílým zákalem \citep{whozakal}, který je relativně snadno a levně léčitelný: Chirurgický zákrok stojí např. v Indii mezi 1 000 a 1 700 USD \citep{cataractsindia}. Z toho vyplývá, že většina nevidomých ve světě je velmi chudá.

\section{Co to je Braillovo písmo}

Braillovo písmo je hmatová abeceda, kterou lze číst a psát bez zraku.

Každé písmeno má dva sloupce zvýšených bodů. Podle toho, kolik bodů je zvýšených a kde jsou mezi nimi mezery, poznáme, které písmeno je napsáno.

Braillovo písmo se liší v každé zemi.  V České republice používáme šestibodové písmo pro tisk a osmibodové písmo pro počítače.  V Německu je písmo osmibodové pro obě použití.

Abychom mohli braillská písmena popsat, jednotlivé body jsou označeny čísly.


\begin{table}[!hb]
\centering
\makebox[0pt][c]{\parbox{1\textwidth}{%
\begin{minipage}[b]{0.3\hsize}\centering
% putStrLn $ utf8ToBrailleBox "1⠁2⠂3⠄4⠈5⠐6⠠7⡀8⢀"
\begin{tabular}{|c|c|}
\hline
1\braillebox{1       }&4\braillebox{   4    }\\
2\braillebox{ 2      }&5\braillebox{    5   }\\
3\braillebox{  3     }&6\braillebox{     6  }\\
7\braillebox{      7 }&8\braillebox{       8}\\
\hline
\end{tabular}

číselné označení bodů

\end{minipage}
\hfill
\begin{minipage}[b]{0.68\hsize}\centering

%%%%%%%%%%%


% putStrLn $ utf8ToBrailleBox "a ⠁ b ⠃ c ⠉ d ⠙ e ⠑ f ⠋ g ⠛ h ⠓ i ⠊ j ⠚ k ⠅ l ⠇ m ⠍ n ⠝ o ⠕ p ⠏ q ⠟ r ⠗ s ⠎ t ⠞ u ⠥ v ⠧ x ⠭ y ⠽ z ⠵ ý ⠯ w ⠷ ž ⠮ ů ⠾ á ⠡ ě ⠣ č ⠩ ď ⠹ š ⠳ ň ⠫  ť ⠳ ó ⠪ ř ⠺ í ⠌ é ⠜ ú ⠬ "
\begin{tabular}{|c|c|c|c|c|c|c|c|}
\hline
a \braillebox{1       }&b \braillebox{12      }&c \braillebox{1  4    }&d \braillebox{1  45   }&e \braillebox{1   5   }&f \braillebox{12 4    }&g \braillebox{12 45   }&h \braillebox{12  5   }\\
\hline
i \braillebox{ 2 4    }&j \braillebox{ 2 45   }&k \braillebox{1 3     }&l \braillebox{123     }&m \braillebox{1 34    }&n \braillebox{1 345   }&o \braillebox{1 3 5   }&p \braillebox{1234    }\\
\hline
q \braillebox{12345   }&r \braillebox{123 5   }&s \braillebox{ 234    }&t \braillebox{ 2345   }&u \braillebox{1 3  6  }&v \braillebox{123  6  }&x \braillebox{1 34 6  }&y \braillebox{1 3456  }\\
\hline
z \braillebox{1 3 56  }&ý \braillebox{1234 6  }&w \braillebox{123 56  }&ž \braillebox{ 234 6  }&ů \braillebox{ 23456  }&á \braillebox{1    6  }&ě \braillebox{12   6  }&č \braillebox{1  4 6  }\\
\hline
ď \braillebox{1  456  }&š \braillebox{12  56  }&ň \braillebox{12 4 6  }&ť \braillebox{12  56  }&ó \braillebox{ 2 4 6  }&ř \braillebox{ 2 456  }&í \braillebox{  34    }&é \braillebox{  345   }\\
\hline
ú \braillebox{  34 6  }\\
\hline
\end{tabular}
% -- https://github.com/timthelion/utf8-to-braillebo://github.com/timthelion/utf8-to-braillebox

česká Braillova abeceda

\end{minipage}%
}}
\end{table}


V Braillově abecedě neexistují velká písmena ani čísla. V šestibodovém Braillově písmu se musí použít sekvence znaků. Velká písmena se značí tak, že se před dané písmeno umístí samotný šestý bod. Velké A je tedy psáno jako
% putStrLn $ utf8ToBrailleBox "⠠⠁"
\braillebox{     6  }\braillebox{1       }
%%%%
. Čísla jsou značena tak, že se před jedno z prvních deseti písmen abecedy umístí speciální znak podobný obrácenému L latinské abecedy
% putStrLn $ utf8ToBrailleBox "⠼"
\braillebox{  3456  }
%%%%
Čísla v šestibodovém Braillově písmu tedy vypadají takto:

% putStrLn $ utf8ToBrailleBox "1⠼⠁2⠼⠃3⠼⠉4⠼⠙5⠼⠑6⠼⠋7⠼⠛8⠼⠓9⠼⠊0⠼⠚"
1\braillebox{  3456  }\braillebox{1       }
2\braillebox{  3456  }\braillebox{12      }
3\braillebox{  3456  }\braillebox{1  4    }
4\braillebox{  3456  }\braillebox{145    }
5\braillebox{  3456  }\braillebox{1   5   }
6\braillebox{  3456  }\braillebox{12 4    }
7\braillebox{  3456  }\braillebox{12 45   }
8\braillebox{  3456  }\braillebox{12  5   }
9\braillebox{  3456  }\braillebox{ 2 4    }
0\braillebox{  3456  }\braillebox{ 2 45   }
%%%%

Osmibodové Braillovo písmo má díky dalším dvěma bodům větší možnosti a není tedy nutné používat sekvence více znaků. K označení velkého písmene stačí současně použít sedmý bod. Velké A je tedy
% putStrLn $ utf8ToBrailleBox "⡁"
\braillebox{1     7 }
%%%%
. K označení čísel se používá osmý bod. 1 je tedy
% putStrLn $ utf8ToBrailleBox "⢁"
\braillebox{1      8}
%%%%
. Pro osmibodové Braillovo písmo neexistuje žádný standard \citep{6dotbraille}.  Protože se používá především na počítači, uživatelé si mohou nastavit vlastní kódování.

Existuje Braillovo kódování pro hudební noty i matematické rovnice.

\section{Jak funguje braillský řádek a jak funguje FCHAD}

\subsection{Braillský řádek}

Braillovo písmo se uplatní přdevším při práci na počítači.  Mnoho nevidomých používá hlasový výstup, aby však mohli pochopit a upravit složitě formátované texty, potřebují braillský řádek.

Braillský řádek je elektromechanické zařízení, které zobrazuje jeden řádek braillského textu.  Běžný braillský řádek zobrazí 40 až 80 znaků a stojí dva až osm tisíc USD.  S vyšším počtem znaků cena obvykle stoupá a naopak.  Braillský řádek, který zobrazí jen 12 znaků, stojí jen tisíc dolarů \citep{perkinsdisplays}.

Dnešní braillské řádky používají piezoelektrickou technologii.  Piezoelektrické krystaly jsou krystaly, jejichž velikost se změní, když jimi projde elektrický proud. Piezoelektrické řádky mají 16 dlouhých krystalů v každém znaku.  Samotné krystaly jsou dost drahé a to se odráží v ceně braillského řádku. To není jediný problém.  Ještě vyšší než cena samotného řádku je cena údržby.  Prach, který vnikne do malého otvoru v řádku, může poškodit piezoelektrický mechanismus.  Člověk, který koupí řádek od americké neziskové společnosti Perkins, má také možnost uzavřít smlouvu o údržbě. Ta však ročně stojí mezi 300 a 500 dolarů USD \citep{perkinsdisplays}.

\subsection{FCHAD}

Už několik let pracuji na novém vynálezu, který jsem pojmenoval FCHAD. FCHAD je jednodušší varianta braillského řádku, která je odolnější a levnější.  Součástky použité při výrobě exempláře, s nímž jsem pracoval při této studii, celkem stály méně než sto dolarů.

FCHAD řádek je jednoznačně horší než peizoelektrický řádek.  Čte se na něm pomaleji a je nutné používat obě ruce.  Jeho hlavní předností je velmi nízká cena.

V apendixu 2 najdete fotografie dvou verzí FCHADu, které byly použity během studie.

FCHAD je zkratka pro \uv{Fast Character Display} tj. rychlý zobrazovač písmen.  FCHAD zobrazuje text po jednom písmenu. To však neznamená, že s ním nelze přečíst delší texty.  Jeho součástí je tzv. selektor, kterým se vybírá písmeno k zobrazení. Obrázek 1 B a C jsou dva různé selektory.  I počítačovou myš lze použít jako selektor.

To, že by bylo levné vyrobit přístroje zobrazující text po jednom písmenu, už napadlo více lidí, ale v literatuře neexistuje pro tento druh řádku název. V této studii tedy nazývám všechny jednoznakové řádky FCHAD.  Existuje mnoho přístrojů typu FCHAD, ale žádné z nich nejsou na trhu.  Informace o některých z nich jsou k dispozici v angličtině na mých webových stránkách. Odkaz je uveden v závěru práce.

\section{Význam pedagogiky v oblasti zařízení pro nevidomé}

Sám nevím, proč tolik FCHADů bylo navrženo, aniž by se dostaly na trh.  Když jsem se zeptal několika lidí, kteří s Braillským písmem pracovali, na jejich názor, často jsem dostal cynickou odpověď, že je výnosnější vyrobit drahý piezo než něco levného.

Mnoho vynálezců pracuje v rámci akademie. Sami nemají možnost založit firmu. Třeba vynalezli nástroje, ale nenašli nikoho, kdo by je chtěl vyrábět. K uvedení nového vynálezu na trh je třeba hodně trpělivosti. Sám na svém vynálezu pracuji již více než rok a ještě nejsem ani v polovině cesty.

Velázquez, který píše o přenosných pomůckách pro nevidomé, má jiný názor na to, proč se tolik zařízení ztratí na cestě.  V závěru vlastní studie píše:
\em
\uv{Přes úsilí badatelů a širokou nabídku dostupných technologií je zájem nevidomých nízký. Audioknihy, braillské řádky, bílá hůl a vodicí pes jsou i nadále jejich nejoblíbenějšími pomůckami.

Jakákoli jiná pomůcka je pro nevidomé velká výzva, přičemž kombinace motivace, spolupráce, optimismu a ochoty/schopnosti učit se novým věcem není samozřejmost.}\em  \citep[str. 15, přeložený z angičtiny]{velazquez2010wearable}%Despite efforts and the great variety of wearable assistive devices available, user acceptance is quite low. Audio books and Braille displays (for those who can read Braille) and the white cane and guide dog will continue to be the most popular reading/travel assistive devices for the blind.
% Acceptance of any other portable or wearable assistive device is always a challenge in blind population. Motivation, cooperation, optimism, willingness/ability to learn or adapt new skills is not a combination that can be taken for granted.

Neochotu nevidomých lidí znám z vlastní zkušenosti. První nevidomý člověk, kterému jsem ukazoval svůj nástroj, sice vůbec nepochopil, jak funguje, ale i tak si byl jistý, že není dobrý. Většina nevidomých naštěstí ochotná se učit je.

Velké procento nových technologií musí čekat, až si na ně lidé zvyknou, než mohou být použité.  Například moje matka stále hledá cestu na papírové mapě, ačkoli by ji na internetu našla několikrát rychleji.  Ještě nedávno hledala kontakty v telefonním seznamu a po telefonu zjišťovala informace, které lze jednoduše najít na internetu.  Technologický vývoj není zpomalován omezenou schopností lidstva vyvíjet nové technologie, ale neochotou či neschopností naučit se je používat.  Pedagogové se nachází uprostřed velkých civilizačních pokroků.  Jejich úkolem je učit i přes neochotu žáka a pomáhat mu orientovat se ve světě, který je ve stavu trvalé změny.

\section{Cíle práce}

Cílem práce je vývoj levné alternativy Braillského řádku na principu Fact Character Display (FCHAD) a dále jeho ověření v praxi na vybraném vzorku nevidomých osob. Budu se zabývat analýzou procesu čtení textového výstupu s podporou tohoto nového neznámého nástroje. Ke splnění zvoleného cíle práce formuluji následující výzkumné otázky:

Jakým způsobem lidé reagují na neznámý nástroj během prvních 45 minut používání?  Jak se novému nástroji přizpůsobují?  Jaké jsou jejich dojmy? Jaká je role pedagoga v procesu podpory učení práce s novým nástrojem?

Zajímavým prvkem studie je to, že FCHAD je pro účastníky úplně neznámý.  Nikdy v životě neslyšeli zkratu FCHAD a nikdy se s podobnými nástroji nesetkali.  Použití FCHADu je složité motoricky a konceptuálně.

Svůj FCHAD jsem vyzkoušel s osmi zrakově postiženými. Účastníci studie na něm četli krátký text a pak mi sdělili své dojmy.  V prvních čtyřech případech mi průběh experimentu zkomplikovaly technické potíže, ale z pěti sezení mám platné výsledky.


\chapter{Teorie}

\section{Pedagogický a psychologický výzkum}

\subsection{Role pedagoga v drilování}

Když jsem si poseděl s těma nevidomým lidmi, rychle jsem se dozvěděl, že nevidomé lidi, které nechápou co mají dělat mají obrovskou schopnost sedět bez pochybu.  Musel jsem je nějak vest k poznání nového nastroje.  Postup, který jsem úplně přirozeně a bez myšlenek zvolil byl postupem drilování.  Když jsem nevidomému člověka říkal aby přečetl na nový nastroj neuměli.  Ale uměli číst první písmeno, které se zobrazuje kdy položili ruce na selektor, a s nějaký pomoci, uměli i přestupovat dal na druhý písmeno.  A tak jsme trénovali.

Někdy jste zkusil trénovat v matice sám doma?  S papírovou knihou těžko jde když nepochopíte zadaný úkolu, vůbec nemate možnost pokračovat.  Pokud rozumíte ale nejste jistý, může Vám pomoc ty "odpovědí v zádě."  Mnoho učebnic je nemají, protože méně motivované žáci by jen koukali v zádě, a nepočítali vůbec.  Nemůžete dat odpovědi v zádě když chcete aby méně motivované žáci něco získali z domácí úkolu.

Je jistý, že mrtvý papír nemůže byt dokonalým prostředkem drilování.  Lepší variant najdeme v počítače.  Počítač muže nám říct jestli jsme počítali správně až potom co zkusíme.  Dále nám dává emoční podporu.  Za každý pět správných odpovědí, hraje se zvuk a animace. Taylor v \em Theory of Learning fo r the Mobile Age\em protestuje, že počítač neumí dozvědět kdy žák něco špatně pochopil a pomáhat mu v pochopení protože neumí sám pochopit\citep{taylor2010theory}. Počítač nemůže Vám pomáhat v případě, že nevíte jak se začínat.

V tom případě, je nutný mít učitele.

Hlavní úkoly učitele v drilování jsou potvrzení odpovědí, vysvětlení nejasností a upřímné lhaní.  To je, \uv{ano Tome, máš to správně} nebo \uv{ne, to nebyl úplně přesný}. \uv{No, musíte nejdřív...} anebo \uv{Jak jste to vyřešil u čísla 5?} a konečně \uv{Děláš to moc dobře, jseš hodně chytrý kluk.}

\subsection{Vliv pedagoga v drilování}

Každý pedagog je jiný a to bude mít vliv na výsledky žáků.  Pedagogické vlastnosti jsou, emoční citlivost, schopnost rychle potvrdit odpovědí, jistotu v potvrzení či vysvětlení, správnost potvrzení nebo vysvětlení.

\subsection{Podstata prvních dojmů a neznámost materiálů}

Pedagog je výjimečně podstatní když jde o neznámy materiál.  Mnoho krát, pokročilý žák se učí líp sám než s pomoci.  Jim to upevňuje samostatnosti a schopnosti rozumět složité materiály.  Žák, který předmět ještě nerozumí naopak sám vůbec neumí dělat pokroky.  Sám neumí ani začít.

Je obecně známý fakt, že první dojmy jsou hodně důležité.  Často jen toho jak se mluví o novém oboru v halu mezi žáci ovlivňuje schopnost žáka to rozumět.

Obory a koncepty jsou často kulturně těžké rozumět.  Tím, že v kultuře se říká, že geometrie je těžká se stává těžká.  Kulturně předurčená složitost muže mít vliv i na moje experiment. Ačkoliv FCHAD je kulturně neznámý nastroj, \uv{nová technologie} je.  Člověk, který má rad novou technologie snad bude mít lehčí čas se učit než člověk, který ji nesnese.

\section{Vybrané teorie a modely učení}

V Průchově \em Přehledu Pedagogiky \em najdeme vymezení dva základní pojmy; vzdělávání a výchova.  Podle Průchy výchova je \uv{záměrné působení na osobnost jedince s cílem dosáhnout změn v různých složkách osobnosti} a vzdělávání je \uv{záměrné a organizované osvojování poznatků, dovedností, postojů, aj., typicky realizovaný prostřednictvím školního vyučování}\citep[str. 16-17]{prucha2009prehled}

Drilování může obsahovat implicitní a explicitní výchovu.  Drilování může implicitně upevňovat kázeň anebo s pomoci pedagoga explicitně potvrdit sebejistotu žáka a sílit jeho důvěr v trpělivosti a práci.  Moje studie by ale neměla mít výchovnou část anebo jestli má to je jenom v tom, že by moje žáci neměli odejít s pocitem, že novou technologie je jen zbyteční ztížení.

Primární účelem drilování je vzdělávací.  Však primární účelem pedagogika v svobodných zemí je vzdělávací a nikoliv výchovný. I když drilování je jeden z základních způsobů vzdělávání pedagogická věda moc o drilování nezajímá. Když jsem našel mnoho studie na křivek učení a drilování všechní byly z psychologických časopisů.

\subsection{Rozdělení učení podle Swifta}

V roce 1903 Edgar James Swift vydal tři studie, které měly projevovat tři druhy učení. Uvedu jeho rozdělení zde, protože přímo vztahuje k přísnému drilování.  Další důvod proč uvedeme Swifta je, že Swift psal před Našem současní módou nelidského \uv{objektivismu}.  Proto najdeme mnoho poznámek v jeho dílu, které možní \uv{nepatři do objektivní vědě}.  Hlavně nás zajímá to, že píše do studie \uv{pedagogické rady a doporučení} kde spojuje výsledky studia s jejích uplatnění v pedagogice. %Pedagogic hints and suggestions

Swiftův druhy učení a související studie jsou:

\em získávání dovednost, naučit se dělat\em : Účastnicí se naučili házet dva míči jednou rukou. % The acquisiton of skill, or learning to do;

\em získávání znalost, trénování souvislost \em : Účastník se učil psát a číst těsnopisem.  % The formation of associations, or the acquisiotn of information

\em trénování inhibice a sebeovládání \em : Účastnici se učili ovládat bezděčné blikání. Účastnici seděli a koukali přes sklenění tabuli jak badatel klepá na sklu dřevěním kladívkem.  Účastnici se měli snažit co nemiň blikat.  % The getting of control, or the formation of inhibitions

\citep[str.201,224,230,231]{swift1903studies}

\subsection{Andersonův(Fittův) stupně získávání dovednosti} %stages of skill acquisition

Základní rozdělení pamětí v psychologie je mezi deklarativní(explicitní) znalosti a procedurální(implicitní) znalost.  Když zapamatujete telefonní číslo můžete ho pamatovat či explicitně anebo implicitně. Explicitní paměť je to, když mohu odpovídat na otázku \uv{jaké je Váš telefonní čislo?} a implicitní paměť je když můžete zadat to číslo do telefonů bez myšlení. Někdy se Vám stálo, že jste zapomněl nějaké telefonní číslo ale ještě uměl ho zadat do telefonů?  Procedurální a deklarativní paměť jsou pevně rozdělení.

Procedurální paměť má nevýhodu, že se nemění.  Když umíte zadat telefonní číslo implicitně, děláte to po každý stejně.  Toho vlastnost procedurálního paměti se říká \em procedurální invariace\em . Abychom se učili něco procedurálně musíme nejdřív explicitně rozumět.  I když procedurální paměť je rozdělení od deklarativní paměť učíme procedurální věci přes deklarativní chápání.  Deklarativní paměť nám umožňuje učit se a vylepšovat naše dovednosti.  Procedurální paměť nám automatizovat dovednosti a činit je bez nepřiměřené námahy\citep[str. 105]{robert1994handbook}.

Andersonův model rozděluje tři stupni procedurální paměť:

\em Kognitivní stadium\em : V toto stadiu žák je vyučování anebo se snaží sám pochopit úkol.  Znalost je explicitní \uv{jak a co}.

\em Asociativní stadium\em : Stadium zlepšení a opravy pochopení.  Jak žák začíná procvičit jeho novou schopnost objevuje, že jeho původní pochopení dovednosti nebyl úplně správní.  Upravuje chyby v pochopení a jeho dovednost se postupně zdokonalí.

\em Autonomní stadium\em : Procedurální znalost.  Už žák osvojil dovednost. Všechno jde teď snadno a implicitně.  Umí činit bez myšlenek. Pracuje rychle a skoro bez chyb\citep{o1987some}.

\subsection{Vlastní model drilování}

Předcházející modely jsou až moc psychologické.  Dávají pozornost přímo na žáka a zapomenou na interakce mezi žáka a učitelem.  Snažím se zde navrhovat model drilování kde jsou uvedené nejenom stadia učení ale i role učitele v každém stadiu.

\em Před znalost \em Člověk má určitý znalost už předtím než začíná cvičit. Například, jestli chci zapamatovat data o dějinný válek, mě pomůže znalost o dějinný umění.  Jestli chci se učit psát na klávesnice mě pomůže znalost abecedy. Pokud chci se učit hrát na klavír asi už vím co to je hudba.  Před-znalost není však stejně jako prerekvizitní znalost.  Některé před-znalosti mohou byt prerekvizitou pro osvojení novou dovednost ale jsou spousta znalosti, které pomůžou ale nejsou nutné znát.  Klavírista snaž se učit hrát na flétnu ale nemusíte umět hrát na klavír abyste se učil hrát na flétnu.  Swift odkazuje na vliv před znalosti v jeho studii o házení míčů, kde studenti, kteří měli zkušenost s basebalem naučili líp nový druh házení než studenti bez zkušenosti. Další studie od Jacqueline Thomas naznačuje, že znalost španělština pomůže studentům, kteří si chtějí naučit franštinu\citep{swift1903studies,thomas1988role}.

Před znalost nemusí byt vždy výhodou. Může taky dojít k konflikt mezi před znalosti a stavějící úkolem. V tohle studie jsem dozvěděl, že někdy účastnicí se snažili používat selektor nesprávně způsobem, které napodobuje správné použití piezoelektrický braillský řádek.  I když znalost podobného jazyka pomůže žákům existuje jev nazvaný jazyková interference kde člověk, který umí dvě jazyky slova někdy spleta, jazyková interference je výraznější tím podobnější jsou ty jazyky\citep{brauer1998stroop}.

{\bf Role učitele}: je určit před znalost žáka, a přizpůsobovat vyučování tak, aby žák zbytečně neopakoval věci, které už zná, aby materiál co nejvíc uvazoval na znalost žáka, a aby materiál odpovídal zájmy a sklony žáka.

\em Chápání úkolu\em : To je kognitivní stadia u Andersona. Drilování funguje jenom když úkol je dostatečně jednoduchý, že chápání jednotlivých kroků úkolu není příliš náročné.  Jestli chci se učit hrát na housle, nemusím mít v hlavě model celého procesu, stačí když mám chytrý učitel, který umí rozdělit složitý dovednost na pochopitelné částky.

Často stadia chápání úkolu se probíhá jako demonstrace.  Nejdřív učitel ukazuje jak se něco děla, potom se snaží žák napodobit.  To jde těžko s nevidomým lidmi.

Další způsob chápání úkolu je přečtení anebo poslouchání návodu.

Třetí možnost je explicitní vedení.  Například, když žák se naučí chodit míč bereme jeho ruku a vedeme správnou cestu přes vzduchu aby věděl \uv{jak to má cítit.}

{\bf Role učitele}: je vykládat úkol jasně a zřetelně ale taky tak aby žák pochopil jakou schopnosti úkol má rozvíjet a k čemu jsou dobré.

Je zvlášť důležitý aby vykládání nenechal dojmy v žákův hlavě, že úkol je příliš složitý, anebo \uv{Jenom pro holky, jenom pro sportovci, jenom pro šprty, atd.}

Swift poznamenal v studie o házení míčů, že způsob který účastnici používaly u začátku učení byl silný indikátor úspěchu. Swiftův dojem z toho je, že dobrý učitel by měl dávat žákům dobrý začátek.  \citep[str. 215]{swift1903studies}  Můj osobní zkušenost nesouhlasí s Swiftem tím, že když jsem se snažil pomáhat účastnici správně položit ruce na FCHADu, často vydřeli správné použity jen par vteřin anebo protestovali, že můj způsob použití FCHADu jim nevyhovuje.
%215 swift "as when the learner is assistet by a good teacher"

\em Obecný schopnost\em a \em Konkretní schopnosti\em :

Můžeme pochopit drilování jako seznám podobných úkolu, který občas opakuje.  Například, když se učíme psát na klávesnice a kopíruji řádek "Táta touží po tequile." Úkol je rozdělení na jednotlivých kroku, v každém krok žák by měl psát jedno písmeno.  Ale písmena se opakují, například písmeno 't' se hodně opakuje.  Obecný schopnost je psát na klávesnice ale konkretní schopnost je psát písmeno 't'. Určité konkretní schopnosti se zdokonalují dřív než obecný schopnost.

{\bf Role učitele}: Učitel může zrychlit zdokonalení obecný schopnost tím, že zjistí, které konkretní schopnosti jsou už zdokonalené.  Učitel pak si vybírá vhodnější úkoly.  Dynamický výběr úkolu je dobře pochopení věda\citep{hintzman1976repetition} a i počítač umí vybírat vhodné úkoly pro žáka\citep{anki}.

\em Dosažené úroveň\em : Kdy žák dostatečně dlouho opakuje drilování už další drilování nezlepšuje výkon.  Každý žák má nějaký přirozený schopnost anebo absolutně dosažitelnou úroveň.  To neznamená, že úroveň u který žák se už nevylepší je jeho stejně jako jeho absolutně dosažitelnou úroveň.  Spousta věci mohou obtěžovat jeho zkutečnou schopnost.  Jeden škodní jev je kognitivní bagáž. Kdy žák často opakuje určitou chybu, často dvě věci spleta, tak může dojít k vytvoření slabý místo. Například, kdy žák se naučí psát na klávesnice, jestli opakovaně píše 'b' kdy chtěl 'v'(dva písmena, která jsou přímo vedle sebe na klávesnice) v budoucnosti žák bude zpomalovat kdy má psát 'v'. Podobní jev je emoční bagáž kde žák se cítí, že část úkolu je těžký, a proto má horší výkon.

{\bf Role učitele}: Učitel může škodit či pomáhat v obraně proti kognitivní a emoční bagáž.  Učitel by měl byt citlivý na to, jak je žák unavený a frustrovaní. Učitel by měl taky byt pečlivý tím jak dává rady.  Jestli dává rady příliš moc, kognitivní bagáž tolik explicitních úvah může spíš zpomalit žáka než mu pomáhat.


\chapter{Výzkum}

\section{Popis studie}

\subsection{První Čast: Čtení Tisku}

Učastnici, nebo možno lepe žáci, v teto studiu nejdřiv poseděli u prazdného stolu.  Neveděli o čem to bude, kromě tomu, že budou vyzoušet nový druh braillský řádek.  FCHAD byl na stole, ale položený do zadu. Tak aby nepřekažel ruky.  Hněd od začatku studii, měli ukol.  Měli číst text v Appendix 3.1.  Četli na hlas a zvuk jsem nahraval.  Text požada souhlas o shromaždění dat, taky vysvětli co to je FCHAD a jak se to použiva. Text byl tísknutí na tuhém kusu braillského papíru velký 26.5x30.5cm.  Maximální delka řádku je 36 písmen.  Braill je šestibodové.  Text byl tisknutí obostraně, a na jedne straně papíru je sloupec dír aby mohl byt navazaný.

Text byl tisknutí v Knihovně a tiskárna pro nevidomé K.E.Macana, která je jediný veřejný tiskařský zavod pro nevidomé v Praze.  Proto předpokladam, že tiskové braill byl velmi standardní.

Když učastník dočte text musí čekat par vkteřinu jak připravím na počitač. Pokud ještě nedali souhlas k zveřejnění informaci jsem požadal o souhlas.  Žadný z mych učastníků odmitly, a nikdo, který jsem pozval k studie odmitl.

\subsection{Otazky}

 Potom, jsem se je zeptal nasledujicí otazek:

Jak jste starý?

Jak dlouho čtete braillské písmo?

Jaký je Váš nejviší dosažený úroveň vzdělání?

Jaký je Váš uroveň zrakové postižený?

Když jsem se je zeptal na věk, jsem vždy objasníl, že nepotřebuji přesný věk.  Jsem poznamenal věky přesné, jenom do pěti let.

Zajimavě, všechní učastnici kromě jeden odpověděli "uplně nevidomý" když jsem je zeptal na jejich uroveň nevidomosti.

WHO rozděli zrakové postižený na 5 kategorii. Nejnížší dvě jsou:

\em "Praktická slepota
zraková ostrost s nejlepší možnou korekcí 1/60 (0,02), 1/50 až světlocit nebo omezení zorného pole do 5 stupňů kolem centrální fixace, i když centrální ostrost není postižena, kategorie zrakového postižení 4

Úplná slepota ztráta zraku zahrnující stavy od naprosté ztráty světlocitu až po zachování světlocitu s chybnou světelnou projekcí, kategorie zrakového postižení 5" \em \citep{sonsklasifikace}

To, že nikdo nehlasil "Praktická slepota" nebo "praktická nevidomost" muže znamenat, že oni otazku nepochopili, že maji odpovědět přesně.  Anebo to muže znamenat, že zkutečně byly všechní uplně nevidomé.

\subsection{Čtení na FCHADu}

Po otazek, jsem přemístil FCHAD aby seděl před učastici.  Musel jsem taky vyzkoušet jestli to funguje správně. To byl první zkušenost pro učastnici s FCHADem.  Slyšeli jak svakne ty body.  Potom, co jsem zjistil, že software je souštění správně jsem řekl učastníkům aby začali číst.  Jenom, že samolřejmě neuměli.  Vidomé lidí, když setkaji s novým nastroje činní. Oni koukaji!  Nevidomé lidí nekoukaji.  Seději, a nic nedělaji.  Koukaji do vzdalenosti, bez pochybu.  Myslím si, že se boji, že kdyby hrabali na stroje, by ho mohli rozbijet.  Není to nerozumní obavání.

Musel jsem něco dělat aby svůj učastnici se snažili. Aby se začali učit.  Jsem je nejdřiv dal ruce na stroje, tak jak by měly byt k spravné čtení.

Vzal jsem je levou rukou, a jsem dal tu levou ruku na zobrazovač(viz Obrazek 1-A). Řekl jsem je, "tady se zobrází to jedno písmeno."  Pak jsem je vzal pravou rukou, a dal jsem tu pravou ruku jejích na kursor selektor(vis Obrazek 1-B,C).

Teď je dobrý čas vysvětlit jak přesně vypada a funguje FCHADy, které byli použití v studie.

\section{Popis FCHADu}

\subsection{První selektor}

První učastnici(1-3) použivali jiný kursor selektor než poslední.  To je kvůli technické zavadům v prvním selektoru.   Technologie selektorů je velmí jednoduchý.  Selektor musí mít dva elektrody.  Když učastník dotyka elektrodem, elektrický proud jde přes tělo z jeden elektrodem k drůhem.  Je jedena katoda, a vice anod.  Podle toho přes, kterou anodou teče proud víme, které písmeno je vybrané(selektované).   První selektor měl 12 anod z napínáčků(vis Obrazek 1-C), a katoda v náramku(Obrazek 3).

Napínáčy jsou nepravidelně umístěné, mezi 1.2 a 2 cm od sebe.  Celý selektor je 19cm od prvního sensoru k poslednému.

Tento system moc dobře nefungoval. Když jsem to zkusil doma, fungoval to bezvadně, ale když to pokusil nějaký chlupatější muž náramní katoda nedostala dobrá kontakt a stroj nevěděl kdy je dotyk a kdy ne.

\subsection{Druhý selektor}

Učastníci 4-6 měli lepší selektor(Vis obrazek 1-B). Nový selektor má dvě řádky elektrody. Horní řádek jsou anody, a dolný řadek jsou katody.  Už proud nemusí projet tělem, ale jenom kuži na prstu.  Tím způsobem, je daleko lepší spoj mezi anodou a katodou.  To ještě nebyl dostatečný.  Jsem musel namoučit prsty učastníků vodou aby dostali spolehlivý spoj.

Další vylepšení je v tom, že sensory jsou v pravidelné umístění 0.5cm od sebe.  Jsou taky daleko blíž k sobě.

\subsection{Zobrazovač}
Další součast FCHADu je zobrazovač(vis Obrazy 1-A a 2).  Zobrazovač, který jsem použil v tehle studie žobrazuje jedno osmibodové písmeno.  Zobrazovaci technologie je zase velmi jednoduchý.  Jde to o cívka, který vybaví tyč aby pohyboval nahoru.  Taková cívka se říka solenoid.  Moje solenoidy vydaji až 200 gramů síly\citep{multicomp}.  To je dost aby jejích pohyb bolelo ale každý bod je jenom 0.45cm vysoký.  Neboli to, když nedříte ruku na zobrazovače silně.  Body nespadnou samy.  Čtenař musí davat ruku na zobrazovače aby správně funguval.

Řádky bodů jsou 2cm od sebe, a sloupci 3cm od sebe(pro učastnici 1-3, byly 5cm od sebe).  Celý zobrazovaci prostor je 6x3cm.  Moje tři delší prsti jsou mezi 7 a 8.5cm dlouhe, a proto, dokažu číst písmena na zobrazovače jenom prstama.  Některé ty ženské učastnicy měly kratší prsti a to nešlo.  Plastový baze Zobrazovač je 5.9cm šíroký(8cm u dřívější učastníci) 9.2cm dlouhý 6cm vysoký.

\subsection{Jak to jde dohromadý}

FCHAD je přípojený na počitač a počitač ma v pamětí tabulku 12ti písmen.  Písmena v tabulce vypulní speciální software, která se jmenuje čtečka obrazovky.  Použival jsem čtečku obrazovky Orca 3.6.3. Orca vyplní tabulku podle toho, jaký text je blízko kursoru(to je, kursor, který je na obrazovce, a ne kursor, který je vybrání kursor selektorem).

Selektorem čtenář vybíra, které ztěch dvanaci písmen bude zobrazené na zobrazovače. Když pohybuje ruku na selektor, čtenař slyší hlasitý svaknutí jak ty solenoidy změneji polohu.

Navrhoval jsem celou elektroniku sroje, programoval jsem ovladač, firmware.  To mě dal možnost nahravat fungování celý stroj.  Když jsem dělal ten výzkum, jsem mohl nahravat lokace prstu uživatela, rychlost čtení, polohu ruky.  Napsal jsem další software, který pomaha v shromaždění dat\footnote{\url{https://github.com/timthelion/anonGraph}}.  Mužete prohlidnout shromaždění data zde\footnote{\url{https://github.com/timthelion/fchad-study}}.

Přesněší detaily nastroje najdete v Appendix I.

\section{Moje žaci}

\subsection{Předběžná studie}

\subsubsection{Žaci 1 a 2}

První "žák" očem budu psát jsem Já Váš batdatel.  Já jsem se učil číst rukama(to je číst Braillova písma) před dvěma roky. Zajimal jsem o tom, protože mivám migreny když čtu.  Nějakou dobu, jsem si řekl: "naučím číst braillovo písmo a už nebudu vůbec číst očima".  Bylo to spíš emoční reakce na frustrace než praktický plan.  Rychle jsem se dozvěděl, že braille vůbec nejde číst rychle a, že braillský řádek je zdaleko přilíš drahý koupít jenom vztekem.  Moj zajem o braillova písma se změnil na hněv o tom jak je všechno pro nevídomé předražený, a snahu oběvit něco levnější.

Umím číst rukama, ale nijak plynule.  Už čtu daleko lepe FCHADem než obyčejný braill.  Je lehčí číst FCHADem, protože písmo je velké.

Když jsem se začal číst na FCHAD četl jsem velmi pomalý.  Ještě jsem neznal braillova písma dokonalý, a často jsem spletl f \braillebox{124}, d \braillebox{145}, h \braillebox{125}, a j \braillebox{245}. Teď je už zvladnu bez problemu.  Dozvěděl jsem, že připojený mezi "psychologický rychlost čtení" a "skuteční rychlost čtení".  Čtení se zda rychlejší, tím mín dělam chyby, a tím vic rozumím text.  Čtení zrychluje daleko pomalejší než psychologický rychlost čtení.  Často, jsem si myslel, že jsem dělal pokrok v svému schopnosti čtení ale když jsem koukal donahravky jsem dozvěděl, že rychlost či nezměnil, anebo jen mírně stoupal.

To neznamena, že jsem se nevylepšil.  Po několik desitek hodin zkušenost na stroje už vidím že moje rychlost soupal od přibližně 0.7 písmen za sekund v 22-11-2012, na 0.8-1.5 písmeno za sekund 22-1-2013.  Plynulost ještě stoupal vic.  Abych počital rychlost čtení písmen za sekund, hledam postupný řádky v graphech lokaci prstu. To jsou místa kde jsem dočetl řádek textu bez tomu, že bych se musel vratit do začatku(protože jsem něco nerozuměl).  Plynulost čtení po celém řádce přeměnil od uplný neplynulost(žádné řádky nečtení ke konce bez chyb) k standardu(většina řádků správně přečtené).

Pro mně, použity selektor je zkoro bez myšlenek, ale poznání znaku na zobrazovač stale někdy trva.  Nastavěl jsem stroj tak, že se použiva selektor pravou rukou, a zobrazovač se čte levou. Jsem si to myslel, protože selektor je podobný jako počitačový myš, a myš se použiva pravou rukou(pro pravaci).  Myslím si, že jsem to nastavěl obraceně.  Nastavený byl stejný pro všechní učastnici studie.

Jestli jste četl titul subsekce vite, že jsem daval "žaci 1 a 2" dohromadý.  To je protože můj FCHAD vůbec nefungoval, když jsem ho ukazal prvnímu učastnikovi. Nevím jakým kozlem to bylo, ale když jsem držel Já katodu v ruce a použival stroj, selektor fungoval bezvadně. Když on to vzal do ruky, začal celý stroj zblaznit.

\subsection{Běžná studie}

\subsubsection{Žak 3}

"Žak" 3 je vysokoškolské vzdělana žena ve věku mezi 45 a 50 let. Už čte braillovo písmo 35 let, to je od začatku školní dochazky.  Pracuje jako vysokoškolskou profesorkou němčiny.  Je uplně nevidomá.

Nehrabala s nastrojem když jsem to nastavil. Zase, FCHAD trpěl technické problemy. Sensory někdy "citilý" dotyky, které v zkutečnosti nebyly.  Zajimavě, i když FCHAD moc dobře nefungoval ona byla schopná na to číst.

Čas na čtení tisknutý text: 93.3 sekundy

Rychlost čtení tisku: 4.36CPS

Správné použití zobrazovači, čtenař by měl položit ruku na zobrazovače rovně, tak aby mohl citít všechní body na řaz.  Třetí učastnik ale použivala sevření prstu aby citila každého tyču zvlašť.  Jelikož to funguje vyborně pro poznání písmena(je to stoprocentně spolehlivý) nejde tím způsobem číst rychle. Já jsem se snažil ji naučit použivat FCHAD správně.  Když jsem ji řikal, že by měla položit ruku rovně a jsem dal jeji ruku na zobrazovače správně řikala "Já nevím, jak ta ruka jen tak leží mně to nevychovuje, Já se musím ty body zvyknout a osahat trouchu jinak." %10:02.6


\subsubsection{Žak 4}

Čtvrtý žak je 20-25 let starý.  Jeho nejvíší uroveň dostažené vzdělávání je maturita.  Je můj spolužak tady na katedře angličtiny. Je uplně nevídomé a taky trpí problemy se sluchem.  Pracuje s braillským písmem od zakladní školy.

Měl velké potiže s přesunem na ty sensory.

\subsubsection{Žak 5}

Paty žak je 20-25 let starý.  Jeho nejvíší uroveň dostažený vzdělávání je maturita. Je uplně nevidomý.  Taký čte braillova písma od zakladní školy.

Kdy začal číst, zase selektor zle choval.  Musel jsem přenastavit prahy.  To mně trval asi deset minut.  To znamenalo, že potom co on četl o stroje, a potom co on teprve dal ruce na stroje, on seděl a čekal.  On měl zdaleko nejlepší vysledky ze všech učastníků.  Možní čekání mu pomohlo.

\subsubsection{Žak 6}

\subsubsection{Žak 7}

\subsubsection{Žak 8}

\subsubsection{Žak 9}

\section{Shrnutí vysledků}


\chapter{Analýza a srovnání}

\section{Obecné dojmy autora}
Čekal jsem, že to bude trvat déle rozumět stroj ale, že účastnici budou schopní rychlejší číst.  Byl jsem překvapení z toho, že účastnici rozuměli základní fungování stroje skoro hned ale četli hodně pomalu.  Některé účastnici skoro nezlepšili během setkání.

Začal jsem studie tím, že jsem chtěl vyzkoumat jak nevidomé lidí samy od sebe objevuji funkčnost nastroje.  Moje účastnici byli až moc pasivní na to, a vždy čekali na můj příkaz.  Asi se báli stroj, který vydá strašlivé cvaknutí když do ně dotykáte.

Měl jsem na mysl, že Andersonův kognitivní stadium bude trvat dlouho, že budou se strojem hrát.  Třeba, kdybych čekal déle by hráli.  Nikdy jsem je nedal přikázaní, že musejí čekat na moje potvrzení správné čtení aby pokračovali. Však, jsem je nikdy neříkal \uv{teď budeme drilovat čtení po jednotlivých písmenech}. Některé zvládly stroj dostatečně dobře, že samostatně četli par písmen ale do větší míru než ne spadli jsme do rigidní drilování, které jsem nikdy explicitně nepřikazoval.  Během předběžní studium, stroj ještě byl dost vadní, že ta rigidnost vychovávala ale u běžní studium už nebyl k tomu důvod.  Asi \uv{čtení} má v sobě určíte míru rigiditu.  Čteme z levá do práva od začátku ke konci není jiný způsob.

Účastnici byly snad jen tak zvykly čekat a poslouchat.

\section{Vymezení fyziologické schopnosti}
Když se zeptáme proč se někdo nějaký schopnost dokázal učit nebo ne musíme se zeptat jestli byly schopní se to naučit.  Protože děláme studium, který zabývá s citem hmatu, musíme najít omezení hmatu.

Abych zajistil, že můj FCHAD je fyzické použitelný jsem si přečetl kapitolu \em Haptic Interfaces\em  od knihy \em HCI Beyond the GUI\em. Člověk má šest druhů hmatových receptorů, volná nervová zakončení, Meissnerovo tělísko, Merkelova buňka, Paciniho tělíska, receptory vlasového folikulu, Ruffiniho tělíska. Různé hmatné receptory mají jiné citlivosti.  Například Paciniho tělíska cítí jenom vibrace.

Pro účelu FCHADu, chceme spustit receptory, které reagují rychle a přesně.  Myslím si, že je to taky lepší spouštět vice druhů receptorů ale nejsem si jistý.  Prostorově nejpřesnější receptory jsou Meissnerovo tělísko a Merkelova buňka.  Ty jsou dost hluboké v kůže, proto používáme silné pohyblivé solenoidy.  Druhý verze FCHADu taky produkuje malou vibrace o 200Hz což je doporučený frekvence v publikace HCI\citep[str. 29-30]{nielsen2008gesture}.

Autoři samy kapitole rozdělí čtyři subjektivní druhy hmatu: tlak, dotyk, vibrace a lechtání.  Myslím si, že lechtání se produkuje pohybem přes prsti.  Jeden hypotéza proč FCHADy nefungují moc dobře je, že pohyb bodů na prstech pomůže čtenář rozpoznat písmeno.  Christophe Ramstein v 1996 studie píše o FCHADech: \em \uv{Protože vice mechanorecoptorů jsou použití v hmatovém vnímání, pohyb prstu nad tvarem anebo zrnitosti zlepšuje vnímání.  V případě jedno-písmenové zobrazovače důsledkem je, že prst necítí přechod mezi písmenem ale jen změna uspořádání.  Protože prst nechybuje, rychlost je pomalejší než s 20ti, 40ti, anebo 80ti písmenová zobrazovače(46.1 slov za minut u rychlé čtenáři a 12.0 slov za minut pro pomalé čtenáři} \em \citep[str. 38, přeložený z angličtiny]{ramstein1996combining}
%"Because multiple mechanical receptors are involved in the tactile perception mechanism, the displacement of the finger over a form or texture improves perception and facilitates acknowledgement[7,14]. As a result, in the single cell case, the finger will not feel the transition from one character to the next, but only the configurational changes in the cell’s pins according to the following character.  Since the finger is not moving on the braille characters, speed performance will of course be slower than on 20-, 40- or 8O-cell braille displays (46.1 wpm for fast readers and 12.0 wpm for slow readers)."
Když si čtu tisknutá Braillova písma se cítím, že je to určitý druh lechtání.  Abych přidal podobnou pocit u zobrazovače FCHADu dal jsem kus látky nad zobrazovačem.  Líbilo mě to ale nepoužíval jsem tu látku v studie. Myslel jsem si, že by věci jenom zkomplikovali pro moje účastnici.

Publikace HCI sestavuje pro naše použity prahy rozlišený anebo v angličtině "Just Noticeable Differences(JND)". Neberu JNDy jako měřítku to co člověk dokáže cítit, ale spíš co nedokáže.  Nechtěl jsem aby nastroj pracoval na hranice citlivosti ale aby znaky byly bezpečně rozpoznatelné.

Pocit hmatu je přesný do ${}1mm^2 - 2mm^2$ na konce prstů a až ${}45mm^2$ jinde na tělo.  Tisknutí Braillovo písmo má body vzdálené od sebe 2.3-2.5mm\citep{brailleautority}. To je až příliš blízko JND.  Však Perkins produkuje speciální velko-písmenové psací stroj kvůli tomu, že některé nevidomé nedokážou normální Braillovo písmo číst\footnote{\url{https://secure2.convio.net/psb/site/Ecommerce?FOLDER=1180&store_id=1101&JServSessionIdr004=9p1uhglug2.app246a&__utma=184864936.426189431.1365936556.1365946389.1366050996.3&__utmb=184864936.1.10.1366050996&__utmc=184864936&__utmx=-&__utmz=184864936.1365936556.1.1.utmcsr=(direct)|utmccn=(direct)|utmcmd=(none)&__utmv=-&__utmk=194396717}}.  Bohužel, když děláme větší písmo už to nevejde na koncem prstu a kůže jinde na tělo je méně citlivá.  Vzdálenost JND kůže na dlaně je 11mm a proto, hrozí se že zvětšení písmo zhoršuje čtení! Můj osobní pocit je, že optimální velikost pro Braillovo písmo je 500 x 300mm.  Když jsem vyrobil první FCHAD jsem nemohl dostat solenoidy tak malé, a proto v studie používal jsem daleko větší zobrazovač\citep[str. 30-32]{nielsen2008gesture}.

Další práh rozlišený je časový.  To není jen jak rychle něco cítíte ale kolik krát můžete něco cítit za sekund.  V angličtině "successiveness limen(SL)" je minimální čas dvě impulsy mohou mít aby člověk je rozlišil od sebe.  V našem studie žádný účastník nečetl tak rychle aby pobližil SL.  Stejně doufal jsem, že budou čist rychlejší, a doufám, že bude možní na FCHADu číst tak rychle jako tisknutí Braillovo písmo.  To je asi 7.9CPS\citep{wetzel2006studies}.  To znamená, že za každý 126ms musí čist jedno písmeno.  SL pro mechanoreceptory je 5 ms ale nás zajímá 20 ms, který je minimální čas aby člověk dvě pocity rozlišil ve správném pořady\citep[str. 32]{nielsen2008gesture}.

Další omezenost je absolutní velikost ruky.  Jestli ruka není dost velká nepokrývá celou zobrazovače.  Několik účastníků měli ruce dost malé aby nepokrývali zobrazovače pohodlně nikdo to nezvládl.

V závěru bych řekl, že účastnici studie byly fyzické schopní používat FCHAD správně a že citové hranici nehráli role v studie.

\section{Podobné nastroje}

\section{Analýza a srovnání s podobnými studii}


\subsection{Swift}
Vrátíme-li k Swiftův rozdělení druhů učení vidíme, že moje studie obsahovala všechní tři druhu.  Účastnicí získali dovednost používat selektor.  Účastnicí trénovali souvislost mezi mezi malé Braillovo písmo, které již znají a velké Braillovo, které se zobrazuje na zobrazovače FCHADu.  Účastnicí odvykli od toho používat selektor jako piezoelektrický braillský řádek a tím pádem trénování inhibice.

\subsection{Anderson}
\subsection{Vlastní model}
Je možný, že to, jak žáci často nemohli říct, který písmo je "Jedná, dva, tři, pět" byla moje chyba jako pedagog jsem nesprávně odhadl před znalost.
\subsection{Optacon}
\subsection{Patterson a Lee}


\chapter{Závěr}
Čekal jsem, že bude žákům trvat déle porozumět stroji.  Také jsem čekal, že budou schopní číst rychleji. Byl jsem překvapen, že většina žáků rozuměla základnímu principu fungování stroje téměř okamžitě. Dalším překvapením bylo, že někteří žáci se během setkání skoro ve čtení nezrychlili.

Zahájil jsem studii s tím, že chci zkoumat, jak nevidomí sami od sebe objevují funkčnost nástroje. Moji žáci byli však příliš pasivní a vždy čekali na můj příkaz. Asi se stroje, který při dotyku vydá strašlivé cvaknutí, báli.

Nikdy jsem jim neřekl, že musejí čekat na moje potvrzení správného přečtení, než budou pokračovat. Nikdy jsem neřekl \uv{teď budeme drilovat čtení po jednotlivých písmenech}. Někteří zvládli stroj natolik dobře, že samostatně četli pár písmen. Do drilování jsme upadli bezděčně, aniž bych ho někdy explicitně vyžadoval. Čekal jsem, že jejich způsob čtení se během setkání zlepší, ale žádnou zásadnější adaptaci jsem nepozoroval. Během předběžné studie byl stroj ještě natolik nedokonalý a zranitelný, že přísné prostředí bylo namístě, během hlavní studie už k němu nebyl důvod. Soudím, že čtení jako takové v sobě obsahuje jistou míru strnulosti. Čteme zleva doprava, od začátku ke konci, není jiný způsob.

Žáci byli zvyklí čekat a poslouchat, neuměli se však adaptovat.

Vrátíme-li se k otázkám uvedeným v kapitole 1.5 Cíle Práce, už můžeme na některé z nich odpovědět.

\em Jakým způsobem lidé reagují na neznámý nástroj během prvních 45 minut používání? Jak se novému nástroji přizpůsobují? Jaké jsou jejich dojmy? \em  Moji účastníci velmi rychle pochopili fungování nástroje, ale nějakou dobu jim trvalo, než byli schopni toto pochopení využít v praxi.  Jejich důvěra ve vlastní schopnost naučit se používat FCHAD přímo odvisela od toho, jestli sami používali braillské řádky.

\em Jaká je role pedagoga v procesu podpory učení práce s novým nástrojem?\em Je jasné, že pro nevidomé je pedagog nezbytnou součástí procesu učení práce s novým nástrojem.  Všichni účastníci potřebovali pomoc s orientací na novém nástroji a opravou chyb při čtení.  Pedagogická intervence byla zvlášť nezbytná u žáka 9, který po většinu doby vůbec nezvládal číst samostatně.

Více informací o této studii naleznete v angličtině na webové stránce: \footnote{\url{https://github.com/timthelion/fchad-study}}.  S veškerými otázkami týkajícími se projektu FCHAD se obracejte na Timothyho Hobbse {\tt timothyhobbs@seznam.cz}.  Technické otázky týkající se návrhů braillských řádků směřujte na emailový seznam
\footnote{\url{ http://mielke.cc/brltty/contact.html#list}}.




\bibliographystyle{csplainnat}
\bibliography{references}
~\\
~\\

\addcontentsline{toc}{chapter}{Bibliography}

\clearpage{}

\appendix
\pagenumbering{Roman}
\thispagestyle{empty}  
\renewcommand{\appendixname}{Appendix}%%přílohy, číslování římskými


\chapter*{Appendix}
\addcontentsline{toc}{chapter}{Appendix}

\section{Technické detaily}

The device comunicates with brltty with the UOBP driver found at \footnote{\url{https://github.com/timthelion/UOBP}}.  The arduino is to be flashed with the FCHAD firmware found at the same url.  The arduino is connected to 8 paralel solenoid circuites as described here \footnote{\url{http://playground.arduino.cc/uploads/Learning/solenoid_driver.pdf}}, using MOSFETs of type IRF530N.  The solenoid power circute is powered by a 6 volt 4 amp wall power supply of type "Smart Electronic Switching adapter Model:JGS1002-24060-2E".  Touches were detected using a 16 channel MUX of type CMOS4067 connected to 4 digital IO pins and one analog pin.  The analog pin was connected to ground with a 10k ohm resistor of type RM0207 "resistor with metalic layer 1\% TK50 0.5W RM 7.5mm".


\clearpage
\begin{figure}
\section{Obrazky FCHADu}
\includegraphics[width=1\textwidth]{Device1}
\begin{subfigure}{.5\textwidth}
  \centering
  \includegraphics[width=.95\linewidth]{Device1-A}
  \caption{Obrazek 1-A - Zobrazovač}
  \label{fig:sub1}
\end{subfigure}%
\begin{subfigure}{.5\textwidth}
  \centering
  \includegraphics[width=.95\linewidth]{Device1-B}
  \caption{Obrazek 1-B - Kursor selektor}
  \label{fig:sub2}
\end{subfigure}

\label{fig:test}
\end{figure}
\begin{figure}
\includegraphics[width=.7\linewidth]{Device1-C}
\caption{Obrazek 1-C - Starý kursor selektor}
\label{fig:sub2}
\end{figure}

\begin{figure}
\centering
\begin{subfigure}{.5\textwidth}
  \centering
  \includegraphics[width=.8\linewidth]{Device2}
  \caption{Obrazek 2 - Zobrazovač ze strany}
  \label{fig:sub1}
\end{subfigure}%
\begin{subfigure}{.5\textwidth}
  \centering
  \includegraphics[width=.8\linewidth]{Device3}
  \caption{Obrazek 3 - Zapesní elektrod}
  \label{fig:sub2}
\end{subfigure}
\label{fig:test}
\end{figure}
\clearpage

%\center
\includegraphics[width=.9\textwidth]{Gymnasium}

Obrazek 4 - Místo konání běžného studia
%\endcenter

\clearpage

\section{Texty}
\subsection{Tisknutý text}
\begin{verbatim}
Děkuji za Vaši účast v této
studii.
Vyzkoušíme nový braillský řádek!

V rámci této studie shromáždím ně-
jaká data o Vás a o tom, jak čtete
Braillovo písmo.  Tato data budou
při zveřejnění anonymní. Jedná se
o Váš přibližný věk, dosažené vzdě-
lání, zkušenost s Braillovým pís-
mem, audionahrávku našeho setkání a
záznam Vaší práce s nástrojem. Než
budeme pokračovat, potřebuji Váš
souhlas se shromážděním a zveřej-
něním těchto dat.

Následující text čtěte prosím na-
hlas a bez přestání až do konce.
Pokud budete mít nějaké otázky,
zeptejte se prosím až potom.

Začátek textu:
Dnes budete seznámen s novým druhem
braillského řádku, který se nazývá
FCHAD.  Tento řádek je navržen
s úmyslem vytvořit nejlevnější možný
                                   1

---- konec stranky ----
\end{verbatim}
\clearpage
\begin{verbatim}

zobrazovač vůbec. Zobrazuje text po
jednotlivých písmenech. Zobrazení
se ovládá pomocí řádku senzorů. Po-
ložením prstu na jeden z dvanácti
kovových bodíků vybíráte, které
z písmen v řádku chcete zobrazit.
Používáte obě ruce: jedna vybírá
písmeno, druhá ho čte.

Konec textu.

Teď Vám položím několik otázek
a poté vyzkoušíme nástroj.
 ---------------------
\end{verbatim}

Text má 138 písmen včetně mezery na první straně a 4 znaků velkých písmen, celkém 142 znaků.

Druhý část má 262 písmen včetně mezery a 3 znaků velkých písmen, celkém 265 znaků.

Celý text má 407 znaků.

\clearpage
\subsection{Text na FCHAD}



Tady je čast textu, který dočetl nejlepší čtenář s braillským překladem.
%1
\begin{tabular}{|c|c|c|c|c|c|c|c|c|c|c|c|}
\hline
B&r&a&i&l&l&s&&&&&\\
\braillebox{1278}&\braillebox{1235}&\braillebox{1}&\braillebox{24}&\braillebox{123}&\braillebox{123}&\braillebox{234}&\braillebox{}&\braillebox{2358}&\braillebox{123}&\braillebox{}&\braillebox{}\\
\hline
\end{tabular}
%2
\begin{tabular}{|c|c|c|c|c|c|c|c|c|c|c|c|}
\hline
k&ý& &ř&á&d&e&k&,& &k&t\\
\braillebox{1378}&\braillebox{12346}&\braillebox{}&\braillebox{2456}&\braillebox{16}&\braillebox{145}&\braillebox{15}&\braillebox{13}&\braillebox{2}&\braillebox{}&\braillebox{13}&\braillebox{2345}\\
\hline
\end{tabular}
%3
\begin{tabular}{|c|c|c|c|c|c|c|c|c|c|c|c|}
\hline
e&r&ý& &t&e&ď& &p&o&u&ž\\
\braillebox{1578}&\braillebox{1235}&\braillebox{12346}&\braillebox{}&\braillebox{2345}&\braillebox{15}&\braillebox{1456}&\braillebox{}&\braillebox{1234}&\braillebox{135}&\braillebox{136}&\braillebox{2346}\\
\hline
\end{tabular}
%4
\begin{tabular}{|c|c|c|c|c|c|c|c|c|c|c|c|}
\hline
í&v&á&t&e&,& &n&e&n&í& \\
\braillebox{3478}&\braillebox{1236}&\braillebox{16}&\braillebox{2345}&\braillebox{15}&\braillebox{2}&\braillebox{}&\braillebox{1345}&\braillebox{15}&\braillebox{2345}&\braillebox{34}&\braillebox{}\\
\hline
\end{tabular}
%5
\begin{tabular}{|c|c|c|c|c|c|c|c|c|c|c|c|}
\hline
p&r&v&n&í& &ř&á&d&e&k&,\\
\braillebox{123478}&\braillebox{1235}&\braillebox{1236}&\braillebox{1345}&\braillebox{34}&\braillebox{}&\braillebox{1235}&\braillebox{16}&\braillebox{145}&\braillebox{15}&\braillebox{13}&\braillebox{2}\\
\hline
\end{tabular}
%6
\begin{tabular}{|c|c|c|c|c|c|c|c|c|c|c|c|}
\hline
 &k&t&e&r&ý& &z&o&b&r&a\\
\braillebox{78}&\braillebox{13}&\braillebox{2345}&\braillebox{15}&\braillebox{1235}&\braillebox{12346}&\braillebox{}&\braillebox{1356}&\braillebox{135}&\braillebox{12}&\braillebox{1235}&\braillebox{1}\\
\hline
\end{tabular}
%7
\begin{tabular}{|c|c|c|c|c|c|c|c|c|c|c|c|}
\hline
z&u&j&e& &j&e&n&o&m& &j\\
\braillebox{135678}&\braillebox{136}&\braillebox{245}&\braillebox{15}&\braillebox{}&\braillebox{245}&\braillebox{15}&\braillebox{1345}&\braillebox{135}&\braillebox{134}&\braillebox{}&\braillebox{245}\\
\hline
\end{tabular}
%8
\begin{tabular}{|c|c|c|c|c|c|c|c|c|c|c|c|}
\hline
e&d&n&o& &p&í&s&m&e&n&o\\
\braillebox{1578}&\braillebox{145}&\braillebox{1345}&\braillebox{135}&\braillebox{}&\braillebox{1234}&\braillebox{24}&\braillebox{234}&\braillebox{134}&\braillebox{15}&\braillebox{1345}&\braillebox{135}\\
\hline
\end{tabular}
%9
\begin{tabular}{|c|c|c|c|c|c|c|c|c|c|c|c|}
\hline
.& & &P&r&v&n&í& &b&y&l\\
\braillebox{378}&\braillebox{}&\braillebox{}&\braillebox{12347}&\braillebox{1235}&\braillebox{1236}&\braillebox{1345}&\braillebox{34}&\braillebox{}&\braillebox{12}&\braillebox{13456}&\braillebox{123}\\
\hline
\end{tabular}
%10
\begin{tabular}{|c|c|c|c|c|c|c|c|c|c|c|c|}
\hline
 &v&y&n&a&l&e&z&e&n& &v\\
\braillebox{78}&\braillebox{1236}&\braillebox{13456}&\braillebox{1345}&\braillebox{1}&\braillebox{123}&\braillebox{15}&\braillebox{1356}&\braillebox{15}&\braillebox{1345}&\braillebox{}&\braillebox{1236}\\
\hline
\end{tabular}
%11
\begin{tabular}{|c|c|c|c|c|c|c|c|c|c|c|c|}
\hline
 &r&o&c&e& &1&9&1&3& &v\\
\braillebox{78}&\braillebox{1235}&\braillebox{135}&\braillebox{14}&\braillebox{15}&\braillebox{}&\braillebox{18}&\braillebox{248}&\braillebox{18}&\braillebox{148}&\braillebox{}&\braillebox{1236}\\
\hline
\end{tabular}
%12
\begin{tabular}{|c|c|c|c|c|c|c|c|c|c|c|c|}
\hline
 &A&n&g&l&i&i&.& & &J&m\\
\braillebox{78}&\braillebox{17}&\braillebox{1345}&\braillebox{1245}&\braillebox{123}&\braillebox{24}&\braillebox{24}&\braillebox{3}&\braillebox{}&\braillebox{}&\braillebox{2457}&\braillebox{134}\\
\hline
\end{tabular}
%13
\begin{tabular}{|c|c|c|c|c|c|c|c|c|c|c|c|}
\hline
e&n&o&v&a&l& &s&e& &O&p\\
\braillebox{1578}&\braillebox{1345}&\braillebox{135}&\braillebox{1236}&\braillebox{1}&\braillebox{123}&\braillebox{}&\braillebox{234}&\braillebox{15}&\braillebox{}&\braillebox{1357}&\braillebox{1234}\\
\hline
\end{tabular}
%14
\begin{tabular}{|c|c|c|c|c|c|c|c|c|c|c|c|}
\hline
t&o&f&o&n& &a& &p&ř&e&v\\
\braillebox{234578}&\braillebox{135}&\braillebox{124}&\braillebox{135}&\braillebox{1345}&\braillebox{}&\braillebox{1}&\braillebox{}&\braillebox{1234}&\braillebox{2456}&\braillebox{15}&\braillebox{1236}\\
\hline
\end{tabular}
%15
\begin{tabular}{|c|c|c|c|c|c|c|c|c|c|c|c|}
\hline
á&d&ě&l& &s&v&ě&t&l&o& \\
\braillebox{1678}&\braillebox{145}&\braillebox{126}&\braillebox{123}&\braillebox{}&\braillebox{234}&\braillebox{1236}&\braillebox{126}&\braillebox{2345}&\braillebox{123}&\braillebox{135}&\braillebox{}\\
\hline
\end{tabular}
%16
\begin{tabular}{|c|c|c|c|c|c|c|c|c|c|c|c|}
\hline
n&a& &z&v&u&k&.& & &K&d\\
\braillebox{134578}&\braillebox{1}&\braillebox{}&\braillebox{1356}&\braillebox{1236}&\braillebox{136}&\braillebox{13}&\braillebox{3}&\braillebox{}&\braillebox{}&\braillebox{137}&\braillebox{145}\\
\hline
\end{tabular}
%17
\begin{tabular}{|c|c|c|c|c|c|c|c|c|c|c|c|}
\hline
y&ž& &č&t&e&n&á&ř& &p&o\\
\braillebox{1345678}&\braillebox{2346}&\braillebox{}&\braillebox{146}&\braillebox{2345}&\braillebox{15}&\braillebox{1345}&\braillebox{16}&\braillebox{2456}&\braillebox{}&\braillebox{1234}&\braillebox{135}\\
\hline
\end{tabular}


\vspace{10mm}


Tady je celý připravený text.

\begin{multicols}{3}
\begin{verbatim}
Braills
ký řádek, kt
erý teď použ
íváte, není 
první řádek,
 který zobra
zuje jenom j
edno písmeno
.  První byl
 vynalezen v
 roce 1913 v
 Anglii.  Jm
enoval se Op
tofon a přev
áděl světlo 
na zvuk.  Kd
yž čtenář po
hyboval spec
iálním perem
 nad písmene
m, přístroj 
bzučel.  To 
bylo velmi p
omalé.  V še
desátých let
ech byl v Ka
lifornii vyn
alezen lepší
 nástroj.  T
akzvaný Opta
con převáděl
 světlo na v
ibrace a zob
razoval pomo
cí 144 drobn
ých tyček.  
To už bylo r
ychlejší, al
e přístroj b
yl hodně dra
hý.  FCHAD, 
který použív
áte právě te
ď, je dosud 
nejlevnější 
hmatový brai
llský zobraz
ovač.  Do bu
doucna by mo
hl být dostu
pný i v chud
ých zemích.
\end{verbatim}
\end{multicols}



\end{document}
